% =======================================================================
% aria_overview.tex
% "Writing the Second Edition with an AI Collaborator"
% 20-minute overview — ARIA 2026
% Author: Martin F. Grace
% Compiled with: pdflatex (MiKTeX, Windows)
% Compile: pdflatex aria_overview  (twice for ToC/nav)
% =======================================================================
\documentclass[11pt, aspectratio=169]{beamer}
\usepackage{tikz}
\usepackage{erm_beamer_theme}
\usepackage{amsmath,amssymb}
\usepackage{booktabs}

% -----------------------------------------------------------------------
% Presentation metadata
% -----------------------------------------------------------------------
\title[AI-Assisted Development of Comprehensive Course
Materials]{Writing the Second Edition of an ERM Text}\\

\author{Martin F.\ Grace}
\institute{%
  University of Iowa\\[2pt]
  Tippie College of Business\\[2pt]
  Vaughan Institute of Risk Management}
\date{ARIA Annual Meeting, 2026}

% -----------------------------------------------------------------------
\begin{document}
% -----------------------------------------------------------------------

% --- Title page ---------------------------------------------------------
\begin{frame}[plain]
\titlepage
\end{frame}

% --- Roadmap ------------------------------------------------------------
\begin{frame}{20-Minute Roadmap}
\begin{enumerate}
  \item \textbf{What the project was} --- 2d Edition of an ERM textbook
  \item \textbf{What did the collaboration entail} --- tools, workflow, division of labor
  \item \textbf{Building a consistent style system} --- \texttt{.sty}, style guides, R themes
  \item \textbf{Chapter-by-chapter production} --- cases, rewrites, what ``engaging'' means
  \item \textbf{From chapters to book} --- integration, cover, index, slides
  \item \textbf{Quality control} --- facts, math, voice
  \item \textbf{What this implies for academic authorship}
  \item \textbf{Open questions for discussion}
\end{enumerate}
\end{frame}

% =======================================================================
\section{The Project}
% =======================================================================
\begin{frame}{What was the Need?}
\begin{frame}{What Triggered the Second Edition}
\begin{itemize}
  \item First edition: standard ERM arc --- foundations, identification, measurement,
        appetite, portfolio, loss control, risk financing, governance.
  \item Problems the first edition left behind:
        \begin{itemize}
          \item Structurally repetitive in places
          \item Visually sparse
          \item More enthusiastic about lists than about \emph{stories}
          \item Two content gaps: financial derivatives, AI as enterprise risk
        \end{itemize}
  \item Audience: juniors/seniors in business programs  
  \item Cconstraints: corporate (not insurance-centric),
        U.S.-oriented. 
\end{itemize}
\end{frame}

\begin{frame}{What Actually Happened (June 3--11, 2026)}
\begin{columns}[T]
\column{0.55\textwidth}
\textbf{Nine chapters, one week:}
\begin{itemize}
  \item Foundations and institutional context
  \item ERM frameworks (COSO, ISO 31000)
  \item Risk identification
  \item Quantification and measurement
  \item Risk appetite
  \item Risk portfolio and correlation
  \item Loss control
  \item Risk finance and ART
  \item Governance
\end{itemize}
\column{0.42\textwidth}
\textbf{Also delivered:}
\begin{itemize}
  \item Master \texttt{erm\_book.tex}
  \item Cover art
  \item Title/copyright/preface
  \item Consolidated bibliography
  \item \texttt{makeindex} index
  \item Nine Beamer slide decks
  \item R figure scripts
\end{itemize}
\end{columns}
\end{frame}

% =======================================================================
\section{Collaboration Architecture}
% =======================================================================

\begin{frame}{Tools and Their Roles}
\begin{block}{The cast}
\begin{itemize}
  \item \textbf{Cursor (Composer agent)} --- project host, integration tasks, slides
  \item \textbf{Claude 4.6 Sonnet / Claude 5 Fable} --- most drafting and \LaTeX{} conversion
  \item \textbf{Perplexity} --- early case research
  \item \textbf{MiKTeX / pdflatex} --- compilation on Windows
\end{itemize}
\end{block}
\medskip
The workflow was a \emph{project}, not a single chat session.
Repository as source of truth; agent transcripts as accidental lab notebook.
\end{frame}

\begin{frame}{The Repeating Chapter Cycle}
\begin{enumerate}
  \item I supply: prior manuscript + intro case draft + pointers to style guide
  \item Model \emph{proposes questions} before large rewrites (enforced after early surprises)
  \item Model produces: revised Markdown chapter, then matching \texttt{.tex}
  \item I compile \texttt{pdflatex}, inspect PDF, send targeted correction requests
        (math, tone, structure)
  \item Style guide updated: anything new (box types, appendix handling, slide theming)
\end{enumerate}
\medskip
\begin{block}{Side branches every cycle}
Update \texttt{.sty}~/ style guide \quad R figures \quad Beamer slides \quad index \& bibliography
\end{block}
\end{frame}

\begin{frame}{Division of Labor}
\begin{center}
\small
\begin{tabular}{@{}p{0.46\linewidth} p{0.46\linewidth}@{}}
\toprule
\textbf{AI assistant} & \textbf{Author} \\
\midrule
First-draft exposition and section scaffolding & Outline, learning objectives, emphasis \\[4pt]
Converting Markdown to consistent \LaTeX{} & Economic reasoning and regulatory nuance \\[4pt]
Applying style environments (boxes, colors) & Choosing real cases and verifying facts \\[4pt]
Generating Beamer slide skeletons & Lecture pacing and institution-specific framing \\[4pt]
Flagging structural repetition & Deciding what to cut \\[4pt]
Patiently regenerating after ``restart'' & Explaining why the third COSO rewrite still reads like a memo \\
\bottomrule
\end{tabular}
\end{center}
\medskip
Pattern: models excel at \textbf{structure} and \textbf{fluency};
humans retain \textbf{judgment} and \textbf{accountability}.
\end{frame}

% =======================================================================
\section{Building a Style System}
% =======================================================================

\begin{frame}{Why a Style System?}
\begin{itemize}
  \item A textbook that looks like nine different books
        teaches students that ERM integration is optional.
  \item Chapter~1 established the visual language:\\
        \alert{Iowa gold and black palette}, learning-objective boxes,
        chapter-overview panels, corporate cases in blue, regulatory notes
        in gray, pitfalls in red.
  \item From Chapter~2 onward: the assistant was asked to \emph{codify} that system.
\end{itemize}
\medskip
\begin{block}{Prompt engineering by documentation}
Instead of re-explaining box conventions in every session, write
``follow \texttt{erm\_textbook\_style\_guide.md}'' --- and get structurally
consistent output.
\end{block}
\end{frame}

\begin{frame}{The Style Infrastructure}
\begin{itemize}
  \item \texttt{erm\_textbook\_style.sty} ---
        \LaTeX{} package: colors, box environments, running heads, chapter-opening layout
  \item \texttt{erm\_textbook\_style\_guide.md} + \texttt{.tex} ---
        parallel human- and machine-readable specifications, updated chapter by chapter
  \item \texttt{erm\_textbook\_style.r} ---
        ggplot theme matching the print palette for figures
  \item \texttt{intro case style guide v2.md} ---
        narrative rules: hook, context, unfolding risk, organizational behavior,
        consequences, conceptual landing; \emph{no} ``as we will see in this chapter''
  \item \texttt{textbook\_style\_manual.md} ---
        voice rules including a prohibition on one-sentence paragraphs
        (``this is a textbook'')
\end{itemize}
\medskip
\small This infrastructure is the part of the project most likely to survive
the next model upgrade.
\end{frame}

% =======================================================================
\section{Chapter Production}
% =======================================================================

\begin{frame}{Opening Cases --- The Hardest Editorial Problem}
\begin{columns}[T]
\column{0.55\textwidth}
\textbf{Cases used:}
\begin{itemize}
  \item Ch.\ 1: J\&J / 1982 Tylenol poisonings
  \item Ch.\ 2: GM ignition switches and siloed risk
  \item Ch.\ 3: Netflix and strategic risk in the DVD era
  \item Ch.\ 5: adidas and risk appetite
  \item Ch.\ 6: Maersk and NotPetya (portfolio correlation)
  \item Ch.\ 9: Wells Fargo, Boeing 737 MAX, governance gaps
\end{itemize}
\column{0.42\textwidth}
\textbf{The problem:}\\
Early drafts ended with\\
``This case illustrates the importance of ERM frameworks,
which will be explored in greater depth\ldots''\\[6pt]
\alert{Technically true. Pedagogically fatal.}\\[6pt]
\textbf{The fix:} Operational rules, not vague instructions.\\
``No second person. No discussion questions. Open with a concrete object.''
\end{columns}
\end{frame}

\begin{frame}{What ``More Engaging'' Actually Meant}
Vague editorial instruction acquired operational meaning over the project:
\begin{itemize}
  \item \textbf{Narrative before taxonomy.}
        Open with a firm in trouble or tension, not with a definition list.
  \item \textbf{Boxes for dense material.}
        COSO, ISO, regulatory notes, and pitfalls break up long exposition.
  \item \textbf{Cross-chapter callbacks.}
        Appetite (Ch.\ 5) echoes identification (Ch.\ 3) and portfolio thinking (Ch.\ 6).
  \item \textbf{Fewer throat-clearing transitions.}
        The AI's default closing move --- ``this will be explored in greater depth'' ---
        was banned.
  \item \textbf{Concrete numbers with caveats.}
        Stylized tax brackets and distress-cost ranges stay, with explicit warnings
        not to treat them as literal policy advice.
\end{itemize}
\end{frame}

\begin{frame}{Notable Chapter Moments}
\begin{itemize}
  \item \textbf{Ch.\ 2 (Frameworks):} COSO boxed, ISO 31000 comparative prose,
        GM case front-loaded. Style guide born after one-sentence-paragraph incident.
  \item \textbf{Ch.\ 4 (Measurement):} Added LTCM/\textit{Moneyball} boxes;
        new non-financial quantification section (sports analytics, supply chain,
        reliability engineering).
  \item \textbf{Ch.\ 7 (Loss control):} Technical appendix from RTF source.
        Two tables failed math review --- ``crazy rounding'' (a phrase I stand behind)
        --- corrected in both Markdown and \LaTeX{}.
  \item \textbf{Ch.\ 8 (Risk finance):} Three legacy part files merged into one chapter:
        captive insurance, catastrophe bonds, ART.
  \item \textbf{Ch.\ 9 (Governance):} Multiple legacy files consolidated;
        Wells Fargo, Boeing 737 MAX, crisis-era gaps aligned with cross-chapter themes.
\end{itemize}
\end{frame}

\begin{frame}{A Representative Timeline (June 3--11)}
\begin{itemize}
  \item \textbf{June 3:} Ch.\ 2 (COSO boxed, GM case) + Ch.\ 3 (identification).
        Wrong chapter number caught. Prohibition on one-sentence paragraphs added.
  \item \textbf{June 4:} Ch.\ 4 with LTCM / \textit{Moneyball} boxes and non-financial
        quantification.
  \item \textbf{June 10--11:} Chapters 5--9 batch-processed. Ch.\ 7 rounding incident.
        Three-part Ch.\ 8 merge. Governance consolidation.
  \item \textbf{Between bursts:} \texttt{pdflatex} Windows path config, MiKTeX setup,
        master \texttt{erm\_book.tex}, cover art, subagent monitoring
        (``are you stuck?'' --- in retrospect, also a governance question).
\end{itemize}
\end{frame}

% =======================================================================
\section{From Chapters to Book}
% =======================================================================

\begin{frame}{Integration: What the Book Build Required}
\begin{itemize}
  \item \texttt{book/erm\_book.tex} --- master file using \texttt{docmute} to include
        standalone chapter preambles without duplication
  \item Cover art (\texttt{erm\_cover\_art.png}), title page, copyright page
  \item Preface: written by me, with honest AI disclosure
  \item Table of contents, consolidated bibliography (\texttt{.bib})
  \item Index (\texttt{makeindex})
\end{itemize}
\medskip
\begin{exampleblock}{What was not automatic}
MiKTeX path configuration, multi-pass compilation, hyperref warnings about
duplicate page anchors on cover pages.
At one point I asked whether the agent was stuck; it was, in fact,
still running subagents. \emph{Patience is a human virtue; context windows are not.}
\end{exampleblock}
\end{frame}

\begin{frame}{The Preface Statement}
\begin{center}
\large\itshape
``This book was written with the help of AI---\\[4pt]
specifically, Perplexity for help with the cases and\\[4pt]
Claude 4.6 Sonnet and 5 Fable for drafting.\\[8pt]
\ldots\ The errors in this book are mine.\\[4pt]
The well-organized paragraphs are at least partly the AI's.''
\end{center}
\medskip
\begin{block}{That division is not false modesty}
It is an accurate model of \alert{accountability}.
\end{block}
\end{frame}

\begin{frame}{Lecture Slides and Figures}
\begin{itemize}
  \item Choice: PowerPoint vs.\ Beamer. We chose Beamer with a custom
        \texttt{erm\_beamer\_theme.sty} matching the textbook palette.
  \item Consistent fonts, gold accents, institute lines for:\\
        University of Iowa $\cdot$ Tippie College of Business $\cdot$
        Vaughan Institute.
  \item Nine chapter decks generated in parallel.
  \item Formatting detail that consumed time: institute line contrast
        (second and third lines rendering near-invisible on black backgrounds).
        Small but real.
  \item Figures from R scripts in \texttt{rscripts/}: J\&J market-share plot,
        aggregate loss distributions, financing spectrum.
        R theme aligns ggplot output with \LaTeX{} colors.
\end{itemize}
\end{frame}

% =======================================================================
\section{Quality Control}
% =======================================================================

\begin{frame}{Facts and Citations --- The Cautionary Tale}
\begin{columns}[T]
\column{0.55\textwidth}
\textbf{The hallucinated citation:}
\begin{itemize}
  \item Model produced: journal, volume, pages, authors.
  \item Everything except reality.
  \item Intro cases constrained to public sources:\\
        10-Ks, Valukas on GM, established LTCM histories,
        major press on Maersk / NotPetya.
\end{itemize}
\column{0.42\textwidth}
\textbf{The fix:}\\[6pt]
Verification, not prompt prettiness.\\[12pt]
\alert{This is also good practice for reading textbooks written by humans.}\\[6pt]
The machine's confidence adds a new flavor of embarrassment
when you almost footnote it.
\end{columns}
\end{frame}

\begin{frame}{Mathematics and Voice}
\begin{columns}[T]
\column{0.5\textwidth}
\textbf{Mathematics:}
\begin{itemize}
  \item Ch.\ 7 appendix: inequalities, present-value expressions,
        worked examples checked.
  \item AI caught some inconsistencies; I caught others reading
        the PDF like a suspicious referee.
  \item Rounding matters: students treat numbers as authoritative
        even when the prose says ``stylized example.''
\end{itemize}
\column{0.47\textwidth}
\textbf{Voice:}
\begin{itemize}
  \item Rules added: no one-sentence paragraphs,
        no AI-ish throat-clearing.
  \item Textbook: second person.
        Intro cases: third-person narrative.
        Style guides enforce the split.
  \item ``landscape of risk'' language: deleted on sight.
  \item Generic examples: replaced with firms and episodes I actually teach.
\end{itemize}
\end{columns}
\end{frame}

% =======================================================================
\section{Implications for Academic Authorship}
% =======================================================================

\begin{frame}{Speed, Scope, and the Productivity Gain}
\begin{itemize}
  \item Nine-chapter revision cycle compressed work that would otherwise have
        spanned many months of evenings.
  \item The AI \emph{lowered the cost of trying structural alternatives}:\\
        box COSO; move the GM case; add a non-financial quantification section.
  \item That is a real productivity gain.
  \item But the bottleneck shifted from \textbf{``produce text''} to
        \textbf{``verify text''} --- with implications for authorial responsibility.
\end{itemize}
\end{frame}

\begin{frame}{Accountability}
\begin{itemize}
  \item Universities have not finished updating authorship norms.
  \item My view: \alert{disclosing AI assistance is necessary};
        treating the model as a co-author is not.
  \item The model has no reputation to lose when the Boeing case dates are wrong.
  \item I remain the author; the AI is infrastructure ---
        sophisticated infrastructure, but infrastructure.
\end{itemize}
\medskip
\begin{block}{A teachable moment for ERM students}
Using AI to write a risk textbook is itself risky: it accelerates drafting
while concentrating quality-control risk in the author.\\
ERM is about deciding which risks to \emph{retain} and which to \emph{transfer}.\\
I retained verification; I transferred first-draft tedium.
\end{block}
\end{frame}

\begin{frame}{What the AI Did Not Reliably Do}
\begin{itemize}
  \item Choose what belongs in a nine-chapter undergraduate treatment versus
        a twelve-chapter aspirational outline
  \item Resist generating plausible nonsense citations
  \item Know Iowa institutional formatting preferences without being told
  \item Finish long integration jobs without monitoring
        (subagents help, but someone must still ask ``are you stuck?'')
\end{itemize}
\medskip
\begin{block}{Future editions}
New content demands: AI risk as a first-class chapter, derivatives coverage expansion,
whatever regulators publish next quarter.
The workflow should transfer if the \textbf{style guides} and
\textbf{repository discipline} persist.
\end{block}
\end{frame}

\begin{frame}{Open Distribution and Reuse}
\begin{itemize}
  \item Repository public under \textbf{CC BY-NC 4.0}.
  \item Instructors can download chapter Markdown, prompt an AI to emphasize
        healthcare or energy, and rebuild.
  \item The README documents compile instructions, directory layout,
        and example customization prompts.
  \item The second edition is therefore not only a text but a
        \alert{template} for AI-assisted course materials ---
        with the caveat that downstream adapters inherit the verification problem.
\end{itemize}
\medskip
\small If the book argues that ERM should be integrated and transparent,
hiding the production stack would be odd.
\end{frame}

\begin{frame}{A Modest Proposal for Departments}
Three policies that would help business schools produce OER textbooks without
pretending the effort is invisible:
\begin{enumerate}
  \item \textbf{Count verified open texts} toward teaching / service portfolios
        when they include reusable assets (slides, code, style packages).
  \item \textbf{Treat AI disclosure} in syllabi and prefaces as standard academic
        practice, not as confession.
  \item \textbf{Fund \emph{verification time}} --- copy editing, fact checking,
        student piloting --- not just drafting time.
\end{enumerate}
\medskip
\begin{alertblock}{The asymmetry}
The marginal cost of \emph{draft prose} has collapsed.\\
The marginal cost of \emph{trusting it} has not.
\end{alertblock}
\end{frame}

% =======================================================================
\section{Conclusion \& Discussion}
% =======================================================================

\begin{frame}{Three Claims Revisited}
\begin{enumerate}
  \item \textbf{Genuine collaboration.} Models did substantial drafting and
        formatting labor. They did \emph{not} set the intellectual agenda.
  \medskip
  \item \textbf{Bottleneck shifted.} From ``produce text'' to ``verify text'' ---
        a shift with real implications for authorial responsibility.
  \medskip
  \item \textbf{Durable infrastructure.} Style guides, \LaTeX{} packages,
        R figure themes, GitHub distribution under CC BY-NC 4.0
        may outlast any single model version.
\end{enumerate}
\end{frame}

\begin{frame}{Discussion Questions}
\begin{enumerate}
  \item How should journals and textbook publishers update disclosure norms?
        Is ``AI-assisted'' a useful category, or does it already describe
        all contemporary academic writing?
  \item The bottleneck has shifted to verification. Should institutions fund
        research assistants specifically for \emph{fact-checking} AI output?
  \item If the marginal cost of a reasonable first draft approaches zero,
        what changes about the economics of open educational resources?
  \item The style guide was the most reusable artifact. What would a discipline-wide
        ERM style guide look like, and who should maintain it?
\end{enumerate}
\end{frame}

\begin{frame}{Thank You}
\begin{center}
{\large Martin F.\ Grace}\\[6pt]
Professor of Finance, Tippie College of Business, University of Iowa\\
Faculty Director, Vaughan Institute of Risk Management\\[10pt]
\texttt{mfgrace@uiowa.edu}\\[14pt]
{\color{ERMDarkGold}\rule{0.5\textwidth}{0.8pt}}\\[10pt]
\small Paper and repository available at:\\
\texttt{https://github.com/ProfMFGrace/erm-book-2ed}\\[6pt]
\textit{Enterprise Risk Management}, Second Edition\\
CC BY-NC 4.0 --- open for customization\\[10pt]
\normalsize\itshape
``Invest early in style guides.
Treat citations as guilty until proven innocent.
Box the boring framework sections.
Write the preface yourself.''
\end{center}
\end{frame}

% -----------------------------------------------------------------------
\end{document}
